\section{The model, normalization, and imported estimates}
\label{sec:model}

We first fix the relation between the two update clocks, the discrete
half-differences, and the heat-bath Dirichlet form.  We then state the
fixed-temperature gap that the proof imports and the two finite-chain
estimates used to control fluctuations and change clocks. We prove the
jump-moment estimate here and the longer clock comparison in
Appendix~\ref{app:clock-comparison}. Both estimates use their stated
finite-chain hypotheses and have no high-temperature restriction.

\subsection{Two clocks and the Dirichlet form}

Let \(\Omega_n=\{-1,1\}^n\).  Independently for \(i<j\), choose
\(J_{ij}\sim N(0,\beta^2/n)\), set \(J_{ji}=J_{ij}\) and \(J_{ii}=0\), and put
\[
 \mu_J(x)=Z_J^{-1}\exp\!\left(\frac12x^{\mathsf T}Jx\right).
\]
The inverse temperature \(\beta\) is fixed as \(n\) tends to infinity.
We suppress \(J\) from \(\mu\) and the transition kernels when this
does not cause confusion.
For \(x\in\Omega_n\), write \(x^i\) for the configuration with spin \(i\)
reversed, and write \(x^{i,+}\) and \(x^{i,-}\) for the configurations
obtained by fixing that spin to \(+1\) and \(-1\), respectively.
Define the half-difference
\[
 d_i f(x)=\frac{f(x^{i,+})-f(x^{i,-})}{2}.
\]
This derivative is independent of \(x_i\).  Writing the conditional
mean, conditional variance, and twice the flip rate as
\[
 m_i(x)=\tanh((Jx)_i),\qquad
 v_i(x)=1-m_i(x)^2,\qquad
 w_i(x)=1-x_i m_i(x),
\]
the generator with one refresh attempt per spin per unit time is
\begin{equation}
 Lf(x)=\sum_{i=1}^n (m_i(x)-x_i)d_i f(x)
      =\sum_{i=1}^n\frac{w_i(x)}2\bigl(f(x^i)-f(x)\bigr).
 \label{eq:model-generator}
\end{equation}
Its semigroup is \(S_t=e^{tL}\).
The kernel \(P\) of one uniformly chosen single-site refresh attempt
therefore satisfies \(L=n(P-I)\).
Both clocks include attempts which leave the configuration unchanged.

We define the unscaled form for any full-support law \(Q\) on the cube
by conditional covariance:
\[
 \mathcal E_Q(f,g)
   =\sum_i \E_Q\Cov_Q(f,g\mid X_{-i}),\qquad
 \mathcal E_Q(f)=\mathcal E_Q(f,f).
\]
When \(Q\) is a Gibbs law with the above interaction and an external
field, let \(v_i\) denote the conditional variance for that law.  Then
\begin{equation}
 \mathcal E_Q(f,g)=\E_Q\sum_i v_i\,d_i f\,d_i g.
 \label{eq:model-dirichlet}
\end{equation}
In particular,
\[
 -\E_\mu[fLg]=\mathcal E_\mu(f,g),\qquad
 \E_\mu[w_i\mid X_{-i}]=v_i,\qquad
 \E_\mu[w_i^2\mid X_{-i}]=v_i.
\]
The last two identities follow by conditioning on \(X_{-i}\) and
averaging a sign with mean \(m_i\).  They explain the two weights:
\(w_i\) occurs in the jump rates and hence in pointwise quadratic
variation, while conditioning converts the stationary energy to the
\(v_i\)-weighted form.  We will make this conversion only inside the
corresponding stationary expectation.

Differentiating the generator later will require the exact product
rule, including its quadratic correction, and the commutation identities
\begin{align}
 d_i(fg)&=f\,d_i g+g\,d_i f-2x_i\,d_i f\,d_i g,
 \label{eq:model-product}\\
 d_i d_jf&=d_jd_if,\qquad d_i^2f=0.
 \label{eq:model-commute}
\end{align}
The zero-field symmetry gives \(\E_\mu X=0\), which will fix the
centering of equilibrium spin observables.  All vector norms are
unnormalized Euclidean norms.
Matrix norms are operator or Hilbert--Schmidt norms unless a normalized
trace \(n^{-1}\Tr\) is written explicitly.
We use \(D_a=\diag(a_1,\ldots,a_n)\).

For a probability density \(g\) relative to \(Q\), write
\[
 \Ent_Q(g)=\E_Q[g\log g],\qquad \E_Qg=1.
\]
For an arbitrary nonnegative \(g\), the homogeneous convention is
\(\Ent_Q(g)=\E_Q[g\log(g/\E_Qg)]\).
Total variation is
\[
 \norm{\nu-\mu}_{\TV}
   =\frac12\sum_{x\in\Omega_n}\abs{\nu(x)-\mu(x)}.
\]
Set
\[
 d_n(t)=\max_{x\in\Omega_n}\norm{S_t(x,\cdot)-\mu}_{\TV},
 \qquad
 d_n^{\rm disc}(k)=\max_{x\in\Omega_n}\norm{P^k(x,\cdot)-\mu}_{\TV}.
\]
Probabilities without a chain subscript refer to the disorder when
the event is a statement about \(J\); conditional path probabilities
are always taken with \(J\) fixed.

\subsection{The spectral-gap input}

The equilibrium input is an all-functions Poincar\'e inequality at a
fixed subcritical temperature.  Its form is the unscaled form just
defined, so its constant controls rate-one-per-site dynamics.

\begin{theorem}[High-temperature gap {\cite[Theorem~1.1]{SKGap}}]
\label{thm:input-gap}
For every fixed \(0<\beta<1\), there is \(\gamma_\beta>0\) such that,
with probability tending to one over \(J\),
\[
 \Var_\mu(f)\leq\gamma_\beta^{-1}\mathcal E_\mu(f)
 \quad\text{for every real function }f\text{ on }\Omega_n.
\]
\end{theorem}

The reference proves this for zero field and Gaussian off-diagonal
variance \(\beta^2/n\), with
\(\mathcal E_\mu=\sum_i\E_\mu\Var(\,\cdot\,\mid X_{-i})\).
Its discrete-time gap is consequently the unscaled gap divided by \(n\).
The same companion supplies deterministic Gaussian-matrix diagnostics
and stable-field residual estimates.  Section~\ref{sec:static} states
those hypotheses before using them.  The arbitrary-law mean estimate in
\Cref{lem:stat-mean} and the simultaneous posterior and cavity gaps in
\Cref{prop:stat-posterior-gap} require additional arguments, which we
supply there.  In particular, \Cref{thm:input-gap} is a zero-field
statement and does not supply a uniform gap at arbitrary external fields.

For any fixed collection of comparisons, intersect the gap event with
the finitely many matrix-diagnostic events those comparisons require.
On the resulting good-disorder event a fixed operator bound
\(\norm J\leq C_\beta\) also holds.
It implies
\begin{equation}
 \min_{x\in\Omega_n}\mu(x)\geq
 \exp\bigl(-(\log2+C_\beta)n\bigr),
 \label{eq:model-minmass}
\end{equation}
because \(\abs{x^{\mathsf T}Jx}/2\leq C_\beta n/2\).
Later arguments selecting countably many comparison tests use a
diagonal sequence of such finite good events; the required
uniformity is stated explicitly in \Cref{thm:path-calculus,prop:cutoff-block}.

\subsection{Two reusable consequences for finite chains}

We now leave the random-matrix input aside.  The first estimate turns a
pointwise gradient bound into exponential moments of a terminal observable.
The second compares averages under shifted Poisson clocks with averages
under a binomial clock.  Both estimates also enter the independent restricted-temperature argument
in the appendices; their proofs use only the finite-chain hypotheses.

\begin{lemma}[Jump exponential moments]
\label{lem:input-mgf}
Fix \(d\geq0\) and \(f:\Omega_n\to\R\).
Suppose a bounded measurable function \(b:[0,d]\to[0,\infty)\) satisfies
\[
 \sup_x\norm{dS_r f(x)}^2\leq b(r)\quad(0\leq r\leq d).
\]
Put \(B=\sup b\) and \(I=\int_0^d b(r)\dd r\).  Then, for every
\(\theta\in\R\) and every starting state \(x\),
\begin{equation}
 \E_x\exp\!\left(\theta(f(X_d)-S_df(x))\right)
 \leq \exp\!\left(2\theta^2 e^{2\abs\theta\sqrt B}I\right).
 \label{eq:model-jump-mgf}
\end{equation}
In particular, \(\Var_x(f(X_d))\leq4I\).
\end{lemma}

\begin{proof}
Use the backward martingale \(F(s,X_s)=S_{d-s}f(X_s)\), whose initial
value is \(S_df(x)\) and terminal value is \(f(X_d)\).  A spin flip
has increment \(-2x_i d_iF\), of magnitude at most
\(2\sqrt{b(d-s)}\).  Its instantaneous sum of squared jump sizes
times rates is
\[
 \sum_i\frac{w_i}{2}(2d_iF)^2\leq4b(d-s).
\]
The backward equation cancels the linear part of the exponential
drift.  The remaining jump compensator is controlled by the preceding
square-jump bound and \(e^z-1-z\leq z^2e^{\abs z}/2\).  Integrating the
resulting drift inequality and dividing by \(e^{\theta S_df(x)}\)
gives \eqref{eq:model-jump-mgf}.  Expansion at \(\theta=0\) gives the
variance assertion.  This argument used the flip-rate bound and the
assumed gradient bound; it used no restriction on the temperature.
\end{proof}

The comparison between lazy discrete clocks and continuous semigroups
follows the shifted Poisson and binomial method of Chen and
Saloff-Coste~\cite[Theorem~1.1 and Lemmas~3.1--3.2]{ChenSaloffCoste2013}.
We give the bounded-array version needed for uniformity in the disorder.
\begin{lemma}[Clock comparison]
\label{lem:input-clock}
Fix $q\in(0,1)$, let $M_n\in\N$ tend to infinity, and let
$(a_{n,j})_{j\ge0}$ be real arrays satisfying
$\sup_{n,j}|a_{n,j}|\le B<\infty$.  Suppose that for every fixed
$y\in\R$,
\begin{equation}
 \E\bigl[a_{n,N_{n,y}}\bigr]\longrightarrow0,
 \qquad
 N_{n,y}\sim\operatorname{Poisson}
                   (q M_n+y\sqrt{M_n}).
 \label{eq:inputclock-poisson-hypothesis}
\end{equation}
The mean in this display is positive for all sufficiently large $n$.
If $B_n\sim\operatorname{Binomial}(M_n,q)$, then
$\E[a_{n,B_n}]\to0$.
The conclusion is uniform over an additional family of arrays and
values of $M_n$ when their bounds and the convergences in
Equation~\eqref{eq:inputclock-poisson-hypothesis} are uniform and $M_n$
tends uniformly to infinity.
\end{lemma}

The proof appears in Appendix~\ref{app:clock-comparison}. It first
approximates the two rescaled clocks by Gaussian densities, then
approximates the narrower Gaussian by a finite signed sum of translates
of the wider one. Fixing that sum before testing the arrays gives the
stated uniformity.


\subsection{The independent case}

The gap input assumes $\beta>0$.  At $\beta=0$, independence gives
the continuous-time cutoff directly, with rate equal to one.

\begin{proposition}
\label{prop:model-independent}
At \(\beta=0\), the continuous-time dynamics has cutoff at
\(\tfrac12\log n\).
\end{proposition}

\begin{proof}
For the upper bound, use the full density relative to equilibrium.
Starting from \(x\), the spins at time \(t\) are independent with
means \(e^{-t}x_i\).  Against uniform measure,
\[
 1+\chi^2(S_t(x,\cdot),\mu)=(1+e^{-2t})^n.
\]
Thus total variation tends to zero uniformly in \(x\) at
\(t=(1+\epsilon)\log n/2\).
For the lower bound, a single statistic suffices.  From the all-plus
state, the normalized sum
\(n^{-1/2}\sum_iX_i\) has mean \(\sqrt n\,e^{-t}\)
and variance at most one; at equilibrium its mean is zero
and its variance is one.
At \(t=(1-\epsilon)\log n/2\), the mean tends to infinity.
Chebyshev's inequality gives total variation tending to one.
\end{proof}

Henceforth \(0<\beta<1\).
The discrete independent case follows from the same clock comparison
used at the end of the paper.
